\documentclass[10pt,a4paper]{amsart}
\usepackage[margin=19mm]{geometry}
\usepackage{amsmath,amssymb,mathtools}
\usepackage[scr=rsfs,cal=euler]{mathalfa}
\usepackage{xfrac}
\usepackage[unicode,hidelinks,bookmarks=false]{hyperref}
\newcommand{\ie}{i.e.\ }
\newcommand{\cf}{cf.\ }
\newcommand{\resp}{respectively\ }
\newcommand{\Subsection}{\subsection}
\providecommand{\nobreakdash}{\nobreak\hbox{-}}
\setlength{\emergencystretch}{2em}
\multlinegap0pt
\usepackage{amssymb}
\usepackage{enumerate}
\usepackage{tikz}
\newtheorem{definition}{Definition}
\newtheorem{remark}{Remark}
\newtheorem{theorem}{Theorem}
\newtheorem{lemma}{Lemma}
\title{Periodic solutions for $p(t)$-Li\'{e}nard equations with a singular nonlinearity of attractive type}
\author{Petru JEBELEAN, Jean MAWHIN, C\u{a}lin \c{S}ERBAN}
\date{Source: arXiv:2506.04927v2 (CC BY 4.0)}
\begin{document}
\maketitle

\noindent\emph{Source and licence.} Periodic solutions for $p(t)$-Li\'{e}nard equations with a singular nonlinearity of attractive type. Version arXiv:2506.04927v2 (CC BY 4.0), distributed by arXiv under the Creative Commons Attribution 4.0 International licence, http://creativecommons.org/licenses/by/4.0/, as recorded in the arXiv licence metadata for this exact version and shown on the abstract page https://arxiv.org/abs/2506.04927v2. Full licence notice: \texttt{licenses/arxiv-2506.04927-CC-BY-4.0.txt}.\par\smallskip

\noindent\textit{Editorial scope.} Counted unit: the article's theorem lemalu (source0.tex lines 433--435). The article's complete original proof of it is delivered with it (source0.tex lines 437--537, one \texttt{proof} environment of 101 lines). No mathematics is written, repaired or re-derived: the statement and the proof are the article's own lines, in the article's own notation, and no retained line is substituted or re-flowed.  Retained uncounted as the source-local context the proof needs: lines 79--86; lines 154--161; lines 174--176; lines 178--179; lines 183--203; lines 425--430.  Nothing was omitted from any retained span, so the package carries no omission record.  No retained line contains a citation, so the wrapper carries no bibliography and no bibliography entry.  No retained line contains a figure environment, an image-inclusion command or drawing code, so no image is delivered and the package declares no asset dependency.  Editorial formatting only, in addition: an AMS \texttt{amsart} preamble that restates the article's own declarations and macros used by the retained lines, namely the environments \texttt{definition}, \texttt{remark}, \texttt{theorem}, \texttt{lemma} on separate counters, together with the standard packages those lines need.  The retained environments print in this document as Definition 1, Remark 1, Theorem 1, Lemma 1 in order of appearance, whereas the article prints the same items with the numbering of its own full document.  The complete native source of the article as distributed in the arXiv e-print for this exact version is bundled unchanged as \texttt{sources/179368/source0.tex}.
\begin{equation}\label{pb1}
	\left \{
            \begin{array}{ll}
		      (\phi_{p(t)}(u'))'+f(u)u'+g(u)=h(t)\quad \mbox{ in }[0,T],\\
		      u(0)-u(T)=0=u'(0)-u'(T),
			\end{array}
\right.
\end{equation}

\begin{equation}\label{pbell}
	\left \{
            \begin{array}{ll}
		      (\phi_{p(t)}(u'))'=\ell(t,u,u')\quad \mbox{ in }[0,T],\\
		      u(0)-u(T)=0=u'(0)-u'(T),
			\end{array}
    \right.
\end{equation}

\begin{definition}\label{def1}\emph{By a {\it solution} of problem \eqref{pbell} we understand a function $u\in C^1_T$ with $\phi_{p(\cdot )}(u'( \cdot ))$ absolutely continuous, which satisfies the differential equation in \eqref{pbell} a.e. on $[0,T]$.
}
\end{definition}

\begin{remark}\label{remark1} \emph{Note that if $\ell:[0,T]\times\mathbb{R}^2\to\mathbb{R}$ is continuous and $u$ is a solution  of \eqref{pbell} in the sense of Definition \ref{def1}, then $\phi_{p(\cdot )}(u'( \cdot ))\in C^1$ and $u$ satisfies the differential equation in \eqref{pbell} everywhere on $[0,T]$ -- this actually means that $u$ is solution in classical sense.}
\end{remark}

\begin{theorem}\label{thHMMT}
Let $\Omega$ be an bounded open set in $C_T^1$ such that the following conditions hold.
\begin{enumerate}
  \item[$(H_1)$] For each $\lambda\in(0,1)$ the problem
  \begin{equation}\label{pblambda}
	\left \{
            \begin{array}{ll}
		      (\phi_{p(t)}(u'))'=\lambda\ell(t,u,u')\quad \mbox{ in }[0,T],\\
		      u(0)-u(T)=0=u'(0)-u'(T)
			\end{array}
  \right.
\end{equation}
has no solution on $\partial\Omega$.
  \item[$(H_2)$] The equation
  $$\mathcal{L}(a):=\frac{1}{T}\int_0^T\ell(t,a,0)dt=0$$
  has no solution on $\partial\Omega\cap\mathbb{R}$.
  \item[$(H_3)$] The Brouwer degree
  $$deg_B[\mathcal{L},\Omega\cap\mathbb{R},0]\neq0.$$
\end{enumerate}
Then problem \eqref{pbell} has a solution in $\overline{\Omega}$.
\end{theorem}

\begin{lemma}\label{elemlema}
Let $a,b\in\mathbb{R}$ with $a<b$ and $\varphi:[0,T]\to\mathbb{R}$ be a continuous function. If there are $\tau_0, \tau_1\in [0,T]$ such that $\varphi(\tau_0)\leq a$ and $\varphi(\tau_1) \geq b$, then there exist $t_0, t_1\in [0,T]$ with $\varphi(t_0)=a$, $\varphi(t_1)=b$ and
$$\mbox{if } t_0<t_1, \mbox{ then }\varphi(t)\in[a,b],\qquad  (t\in[t_0,t_1]),$$
respectively,
$$\mbox{if } t_1<t_0, \mbox{ then }\varphi(t)\in[a,b],\qquad  (t\in[t_1,t_0]).$$
\end{lemma}

\begin{theorem}\label{lemalu}
If problem \eqref{pb1} has a lower solution $\alpha$ and an upper solution $\beta$ such that $\alpha(t)\leq\beta(t)$ for all $t\in[0,T]$,  then it has a solution $u$ with $\alpha(t)\leq u(t)\leq\beta(t)$ for all $t\in[0,T]$.
\end{theorem}

\begin{proof} Let $\gamma:[0,T]\times\mathbb{R}\to(0,\infty)$ be the continuous function given by
$$\gamma(t,x):=\left \{
            \begin{array}{lll}
		      \beta(t),\quad &\mbox{ if}& x>\beta(t),\\
		      x,\quad\ &\mbox{ if}& \alpha(t)\leq x\leq\beta(t),\\
              \alpha(t),\quad &\mbox{ if}& x<\alpha(t)
			\end{array}
\right.$$
and define $\ell^{*}:[0,T]\times\mathbb{R}^2\to\mathbb{R}$ by $$\ell^*(t,x,y)=h(t)-g(\gamma(t,x))-f(\gamma(t,x))y+x-\gamma(t,x).$$ Note that, as $\ell^{*}$ is continuous, below the notion of solution will be understood in classical sense (see Remark \ref{remark1}).  We introduce the modified problem
\begin{equation}\label{pbmod}
	\left \{
            \begin{array}{ll}
		      (\phi_{p(t)}(u'))'=\ell^*(t,u,u')\quad \mbox{ in }[0,T],\\
		      u(0)-u(T)=0=u'(0)-u'(T)
			\end{array}
\right.
\end{equation}
and, in order to apply Theorem \ref{thHMMT}, for $\lambda\in(0,1)$, we consider the homotopy problem associated with \eqref{pbmod}
\begin{equation}\label{pbmodlambda}
	\left \{
            \begin{array}{ll}
		      (\phi_{p(t)}(u'))'=\lambda\ell^*(t,u,u')\quad \mbox{ in }[0,T],\\
		      u(0)-u(T)=0=u'(0)-u'(T).
			\end{array}
\right.
\end{equation}

Next, let $u \equiv u_{\lambda}$ be a solution of problem \eqref{pbmodlambda} for some $\lambda\in(0,1)$. We choose $M_1>\|\beta\|_\infty$ a constant such that
\begin{equation}\label{M1}
    h(t)-g(\beta(t))+M_1-\beta(t)>0\ \mbox{ and }\ h(t)-g(\alpha(t))-M_1-\alpha(t)<0 \quad (t\in[0,T])
\end{equation}
and we show that $\|u\|_\infty<M_1$. Otherwise, suppose that there exists $t_0\in[0,T]$ such that $u(t_0)=\max_{t\in[0,T]}u(t)\geq M_1$ or $u(t_0)=\min_{t\in[0,T]}u(t)\leq -M_1$.
If $t_0\in(0,T)$ and $u(t_0)\geq M_1$, then $u'(t_0)=0$ and since $u$ is a solution of \eqref{pbmodlambda}, on account of $M_1>\|\beta\|_\infty$ and \eqref{M1}, one has
\begin{eqnarray*}
  (\phi_{p(t)}(u'))'_{t=t_0} &=& \lambda\ell^*(t_0,u(t_0),0)=\lambda(h(t_0)-g(\beta(t_0))+u(t_0)-\beta(t_0)) \\
  &\geq& \lambda(h(t_0)-g(\beta(t_0))+M_1-\beta(t_0))>0.
\end{eqnarray*}
So, there exists a positive constant $\delta$ such that $\phi_{p(t)}(u'(t))$ is strictly increasing on $(t_0,t_0+\delta)$. Thus, $$\phi_{p(t)}(u'(t))>\phi_{p(t_0)}(u'(t_0))=0,\quad (t\in(t_0,t_0+\delta)),$$ which shows that $u'(t)>0,\ t\in (t_0,t_0+\delta)$. Hence, $u$ is strictly increasing on $(t_0,t_0+\delta)$, contradicting $u(t_0)=\max_{t\in[0,T]}u(t)$. In a similar way, if $t_0\in(0,T)$ and $u(t_0)\leq-M_1$, since $-M_1<\alpha(t)$, for all $t\in[0,T]$ and using the second inequality in \eqref{M1}, we obtain
$$(\phi_{p(t)}(u'))'_{t=t_0} = \lambda\ell^*(t_0,u(t_0),0)\leq\lambda(h(t_0)-g(\alpha(t_0))-M_1-\alpha(t_0))<0,$$
which implies that there is $\delta>0$ such that $u$ is strictly decreasing on $(t_0,t_0+\delta)$, i.e. again a contradiction with $u(t_0)=\min_{t\in[0,T]}u(t)$.

If $t_0=0$ or $t_0=T$ and $u(0)=u(T) \geq M_1$, then $u'(0) \leq 0$ and $u'(T) \geq 0$ implying that $u'(0)=0$ and arguing as above we get a contradiction. The case when $t_0=0$ or $t_0=T$ under the assumption $u(0)=u(T) \leq -M_1$ can be treated by similar arguments.

Next, we claim that there exists a constant $M_2>0$, such that $\|u'\|_\infty<M_2$. To prove this, observe first that
\begin{equation}\label{nagumo}
    |(\phi_{p(t)}(u'))'|\leq|\ell^*(t,u,u')|\leq\|h\|_{\infty}+c_1+c_2|u'|+2M_1\leq c_3(1+|u'|),\quad (t\in[0,T]),
\end{equation}
where $c_1$, $c_2$ and $c_3$ are positive constants depending on $M_1$. Since $u(0)=u(T)$, there is some $\tau_0\in[0,T]$ such that $u'(\tau_0)=0$. Then, obviously one has
\begin{equation}\label{utau0}
    0=|u'(\tau_0)|^{p(\tau_0)-1}<\left(\frac{2M_1}{T}+1\right)^{p^+-1}.
\end{equation}
\noindent We choose the constant $M_2>1$ such that $$\frac{M_2^{p_--1}-\left(\displaystyle\frac{2M_1}{T}+1\right)^{p^+-1}}{c_3}>T+2M_1$$
and suppose by contradiction that there is a point $\tau_1\in [0,T]$ such that $|u'(\tau_1)|\geq M_2$, implying that
\begin{equation}\label{utau1}
    |u'(\tau_1)|^{p(\tau_1)-1}\geq M_2^{p_--1}.
\end{equation}
In view of \eqref{utau0}, \eqref{utau1} and Lemma \ref{elemlema}, there exist $t_0,t_1\in[0,T]$ with $$|u'(t_0)|^{p(t_0)-1}=\left(\displaystyle\frac{2M_1}{T}+1\right)^{p^+-1},\ \ |u'(t_1)|^{p(t_1)-1}=M_2^{p_--1}$$ and
$$\mbox{if } t_0<t_1, \mbox{ then } |u'(t)|^{p(t)-1}\in\left[\left(\displaystyle\frac{2M_1}{T}+1\right)^{p^+-1},M_2^{p_--1}\right],\ \ (t\in[t_0,t_1]),$$
respectively,
$$\mbox{if } t_1<t_0, \mbox{ then } |u'(t)|^{p(t)-1}\in\left[\left(\displaystyle\frac{2M_1}{T}+1\right)^{p^+-1},M_2^{p_--1}\right],\ \  (t\in[t_1,t_0]).$$
We discuss the case $t_0<t_1$; similar arguments work when $t_1<t_0$. Using the first change of variables formula, we have
\begin{eqnarray}\label{estim0}
  T+2M_1 &<& \frac{M_2^{p_--1}-\left(\displaystyle\frac{2M_1}{T}+1\right)^{p^+-1}}{c_3} =\int\limits_{\left(\frac{2M_1}{T}+1\right)^{p^+-1}}^{M_2^{p_--1}}\frac{1}{c_3}\ ds\nonumber\\
  &=& \int\limits_{|u'(t_0)|^{p(t_0)-1}}^{|u'(t_1)|^{p(t_1)-1}}\frac{1}{c_3}\ ds=\int_{t_0}^{t_1}\frac{(|u'(t)|^{p(t)-1})'}{c_3}\ dt.
\end{eqnarray}
Note that $u'$ does not vanish in $[t_0,t_1]$. If $u'(t)$ is strictly positive on $[t_0,t_1]$, then on account of \eqref{estim0} and \eqref{nagumo}, we obtain
\begin{eqnarray*}
  T+2M_1 &<& \int_{t_0}^{t_1}\frac{(\phi_{p(t)}(u'))'}{c_3}\ dt=\int_{t_0}^{t_1}\frac{\lambda\ell^*(t,u,u')}{c_3}\ dt \\
  &\leq& \int_{t_0}^{t_1}(1+u'(t)) dt\leq T+|u(t_1)-u(t_0)|<T+2M_1,
\end{eqnarray*}
i.e., a contradiction. Similarly, if $u'(t)$ is strictly negative on $[t_0,t_1]$, one gets
\begin{eqnarray*}
  T+2M_1 &<& \int_{t_0}^{t_1}\frac{-(\phi_{p(t)}(u'))'}{c_3}\ dt=\int_{t_0}^{t_1}\frac{-\lambda\ell^*(t,u,u')}{c_3}\ dt \\
  &\leq& \int_{t_0}^{t_1}(1-u'(t)) dt\leq T+|u(t_0)-u(t_1)|<T+2M_1,
\end{eqnarray*}
 a contradiction, again.
\smallskip

Hence, $\|u\|_1<M_0$, where $M_0:=M_1+M_2$, and we shall apply Theorem \ref{thHMMT} with $\Omega:=B(0,M_0)$ and $\ell=\ell^*$.
Clearly,  $(H_1)$ in Theorem \ref{thHMMT} is satisfied.
From \eqref{M1}, we have that $\ell^*(t,-M_0,0)<0$ and $\ell^*(t,M_0,0)>0$ for all $t\in [0,T]$, which yield
$$\mathcal{L}(-M_0)=\frac{1}{T}\int_0^T\ell^*(t,-M_0,0)dt<0,\ \ \ \mathcal{L}(M_0)=\frac{1}{T}\int_0^T\ell^*(t,M_0,0)dt>0$$
and, as $\partial \Omega \cap \mathbb{R}=\{ -M_0, M_0 \}$, hypothesis $(H_2)$  is fulfilled. Finally,  $deg_B[\mathcal{L},\Omega\cap\mathbb{R},0]=1$, hence $(H_3)$ also holds true. Therefore, problem \eqref{pbmod} has at least one solution in $\overline{\Omega}$.

We conclude the proof by showing that if $u$ is a solution of \eqref{pbmod}, then $\alpha(t)\leq u(t)\leq\beta(t)$ for all $t\in[0,T]$. Suppose by contradiction that there is some $t_0\in[0,T]$ such that $$\max_{t\in[0,T]}(\alpha(t)-u(t))=\alpha(t_0)-u(t_0)>0.$$
If $t_0\in(0,T)$, then $\alpha'(t_0)=u'(t_0)$ and there is a sequence $\{t_k\}$ in $(t_0, T)$ converging to $t_0$ such that $\alpha'(t_k)-u'(t_k)\leq0$. As $\phi_{p(t)}: \mathbb{R} \to \mathbb{R}$
is an increasing homeomorphism for all $t\in [0,T])$, this yields
$$\phi_{p(t_k)}(\alpha'(t_k))-\phi_{p(t_0)}(\alpha'(t_0))\leq\phi_{p(t_k)}(u'(t_k))-\phi_{p(t_0)}(u'(t_0)),$$
implying that
$$\left(\phi_{p(t)}(\alpha'(t))\right)_{t=t_0}'\leq\left(\phi_{p(t)}(u'(t))\right)_{t=t_0}'.$$
Thus, because $\alpha'(t_0)=u'(t_0)$ and $\alpha$ is a lower solution of \eqref{pb1}, we obtain
\begin{eqnarray*}
  \left(\phi_{p(t)}(\alpha'(t))\right)_{t=t_0}' &\leq& \left(\phi_{p(t)}(u'(t))\right)_{t=t_0}'=\ell^*(t_0,u(t_0),u'(t_0)) \\
  &=& h(t_0)-g(\alpha(t_0))-f(\alpha(t_0))\alpha'(t_0)+u(t_0)-\alpha(t_0)\\
  &<& h(t_0)-g(\alpha(t_0))-f(\alpha(t_0))\alpha'(t_0)\leq \left(\phi_{p(t)}(\alpha'(t))\right)_{t=t_0}',
\end{eqnarray*}
a contradiction. If $$\max_{t\in[0,T]}(\alpha(t)-u(t))=\alpha(0)-u(0)=\alpha(T)-u(T)>0,$$
then $\alpha'(0)-u'(0)\leq0$ and $\alpha'(T)-u'(T)\geq0$. Since $\alpha$ is a lower solution of \eqref{pb1}, one has $\alpha'(0)\geq\alpha'(T)$ and using this, we get that $\alpha'(0)-u'(0)=0=\alpha'(T)-u'(T)$. It follows  $\phi_{p(0)}(\alpha'(0))=\phi_{p(0)}(u'(0))$. On the other hand, $\max_{t\in[0,T]}(\alpha(t)-u(t))=\alpha(0)-u(0)>0$ implies, reasoning as above, that
$$\left(\phi_{p(t)}(\alpha'(t))\right)_{t=0}'\leq\left(\phi_{p(t)}(u'(t))\right)_{t=0}'.$$
Using this inequality, together with $\alpha'(0)=u'(0)$ and the fact that $\alpha$ is a lower solution of \eqref{pb1}, we again proceed as above to obtain a contradiction. Consequently, $\alpha(t)\leq u(t)$ for all $t\in[0,T]$. Analogously, using that $\beta$ is an upper solution of problem \eqref{pb1}, it follows that $u(t)\leq\beta(t)$ for all $t\in[0,T]$. The proof is now complete.
\end{proof}

\end{document}
